\section{Resultados preliminares}
\label{sec:resultados}

Esta seção apresentará separadamente as evidências do experimento de viabilidade e os resultados do experimento principal. A primeira etapa informa a exequibilidade técnica e os ajustes do protocolo; seus resultados não serão tratados como teste das hipóteses. A análise principal será reportada após a consolidação do conjunto de referência, a execução controlada dos recuperadores e a avaliação das respostas conforme o protocolo.

\subsection{Experimento de viabilidade: execução e ajustes do protocolo}

Os artefatos disponíveis registram quatro execuções completas, todas sem erro de processamento: duas rodadas com 20 perguntas e duas identificadas como piloto de 32. A Tabela~\ref{tab:resultados-viabilidade} resume as médias atribuídas pelo juiz automatizado. As duas rodadas de 20 perguntas usaram arquivos de pergunta distintos. As duas rodadas de 32 apontam para o mesmo nome de arquivo de piloto, mas registram hashes diferentes tanto para a bateria quanto para a rubrica; por isso, não é possível confirmar que os itens eram idênticos, nem tratá-las como repetições estritamente idênticas.

\begin{table}[htbp]
\centering
\caption{Execuções automatizadas registradas no experimento de viabilidade.}\label{tab:resultados-viabilidade}
\begin{tabular}{lrrr}\toprule
Execução (UTC) & Perguntas concluídas & Erros & Média do juiz (0--10)\\\midrule
05/05/2026, 14:26 & 20 & 0 & 9,350\\
05/05/2026, 15:33 & 20 & 0 & 9,250\\
05/05/2026, 16:27 & 32 & 0 & 9,031\\
05/05/2026, 17:11 & 32 & 0 & 9,094\\\bottomrule
\end{tabular}
\par\smallskip Fonte: elaborado pelo autor a partir dos resumos de execução do projeto de viabilidade (2026).
\end{table}

As médias são notas produzidas pela rubrica exploratória do projeto, composta por cinco critérios em escala de 0 a 2, e não métricas de recuperação ou avaliações humanas independentes. Nas quatro configurações registradas, `gpt-5-nano` foi empregado tanto na geração quanto no julgamento. As notas, portanto, descrevem a saída de um avaliador automatizado que compartilha o modelo com o sistema gerador e não devem ser interpretadas como estimativa autônoma da qualidade das respostas.

Os resumos também registram risco alto de mistura de autores entre os trechos recuperados em 29 dos 32 itens de uma rodada e em 28 dos 32 itens da outra. Esse indicador é um sinal para inspeção dos contextos e das atribuições documentais; não equivale, por si só, a erro confirmado em todas as respostas. Na série de abril, uma inspeção proposta para sete respostas apontou que a resposta `q20`, apesar de receber nota automática 10/10, extrapolava o escopo da pergunta. Como essa revisão foi preparada como proposta argumentada e não como anotação independente, o caso ilustra uma hipótese de falha a verificar, sem fornecer medida de concordância humano--juiz.

Uma execução complementar de 5 de maio processou duas perguntas selecionadas para verificar problemas de escopo e recuperação; ambas foram concluídas, com média automática 9,0, e o resumo indicou risco alto de mistura de autores nos dois casos. Pelo tamanho e pela seleção dirigida, essa checagem funciona apenas como diagnóstico pontual e não foi agregada às quatro rodadas completas.

Em conjunto, os registros demonstram que o fluxo de preparação, consulta e avaliação automatizada pôde ser executado com as baterias utilizadas e produzir artefatos rastreáveis. Eles também identificam questões para o protocolo seguinte: controlar autoria dos trechos recuperados, verificar o escopo das respostas e preservar hashes e configurações por rodada. Não demonstram desempenho validado nem testam as hipóteses da dissertação.

\subsection{Caracterização da bateria e do conjunto de referência}

O material de trabalho disponível contém uma bateria de 200 perguntas em 16 categorias e um piloto estratificado de 32 itens, com dois itens por categoria. As categorias abrangem perguntas factuais, foco e desambiguação de autoria, consultas multifocais, verificação de citações, comparação, limites de contexto, confusão entre autores, armadilhas de alucinação, múltiplas etapas, evidência negativa, questões numéricas, estresse de recuperação, controle de escopo e síntese. Essa seleção busca cobrir tipos distintos de consulta, mas não é amostra probabilística de perguntas reais.

A classificação de respondibilidade nos arquivos distingue itens respondíveis, parcialmente respondíveis e não respondíveis. A curadoria do conjunto de referência da dissertação ainda precisa confirmar essas classificações e associar, para cada item respondível, os documentos e trechos suficientes; as fontes indicativas da bateria não constituem julgamentos exaustivos de relevância. Consequentemente, ainda não há estatísticas de um conjunto-ouro validado que permitam reportar resultados inferenciais de recuperação.

\subsection{Resultados da recuperação documental}

Os registros examinados não contêm uma comparação controlada entre BM25, recuperação vetorial e recuperação híbrida com consultas e julgamentos de relevância validados. Assim, não são reportados valores de Recall@k, Precision@k, MRR, nDCG ou Hit@k para o experimento principal. O piloto lexical registrado em outro artefato usa consultas provisórias e não substitui essa comparação. O cálculo das métricas previstas dependerá da validação dos documentos e trechos relevantes e da execução dos recuperadores sob condições equivalentes.

\subsection{Resultados da geração de respostas}

As notas resumidas acima dizem respeito à rubrica exploratória do protótipo, com cinco critérios e escala de 0 a 2 por critério, agregados em escala de 0 a 10. A rubrica de cinco dimensões definida para o experimento principal da dissertação possui finalidade e regras próprias; não há resultados disponíveis de sua aplicação às respostas da bateria curada. Por essa razão, não se apresentam médias por dimensão nem pontuação global da rubrica principal.

\subsection{Avaliação humana e avaliação automatizada}

Os artefatos disponíveis não documentam uma rodada de avaliação humana independente, nem informam perfil, número de avaliadores ou concordância calculada. A inspeção de sete respostas mencionada acima foi uma análise amostral preparada como proposta, sem independência de avaliador e sem cálculo de concordância. Também existe um resumo de estabilidade de abril: três rodadas de 20 perguntas na mesma base tiveram médias automatizadas de 9,15, 9,55 e 9,35, e 12 perguntas variaram entre as execuções. Esse registro descreve variação na configuração e no conjunto examinados; não mede estabilidade geral, e não supre a comparação com julgamentos humanos.

\subsection{Análise de erros e casos críticos}

Os casos documentados orientam uma agenda de inspeção qualitativa, sem permitir estimar a frequência geral de cada falha. A mistura de autores assinalada nos resumos deve ser confrontada com os trechos e metadados de cada resposta para distinguir contexto heterogêneo de atribuição efetivamente incorreta. A inspeção de abril sobre `q20` requer revisão independente para confirmar o diagnóstico de extrapolação de escopo. As fontes examinadas não oferecem uma análise sistemática de casos de insuficiência documental ou de perguntas não respondíveis; esses padrões deverão ser registrados durante a curadoria e a execução do experimento principal.
